% !TeX program = pdflatex
\documentclass[11pt,a4paper]{article}

\usepackage[margin=2.5cm]{geometry}
\usepackage[T1]{fontenc}
\usepackage[utf8]{inputenc}
\usepackage{lmodern}
\usepackage{microtype}
\usepackage{booktabs}
\usepackage{tabularx}
\usepackage{enumitem}
\usepackage{xcolor}
\usepackage[colorlinks=true,linkcolor=blue,urlcolor=blue,citecolor=blue]{hyperref}

\newcolumntype{Y}{>{\raggedright\arraybackslash}X}

\title{\textbf{Explainability Layer}}
\author{gommage}
\date{June 2026}

\begin{document}
\maketitle

\begin{abstract}
gommage is not merely an application with an explainability feature. The whole
project is the explainability layer for agent execution. It records what an
agent saw, what it inferred, which tools it called, which evidence supported
each step, and how the execution unfolded. The central artifact is the
\emph{Agent Execution Record} (AER), which stores reasoning provenance as a
structured object and enables the execution to be visualized as a directed
acyclic graph (DAG).
\end{abstract}

\section{Purpose}

The purpose of the explainability layer is to make autonomous agent behavior
inspectable after execution. Traditional logs show events, but they rarely show
the reasoning context that caused an event. gommage captures the agent's
trajectory as structured provenance so developers, auditors, and operators can
answer three questions:

\begin{enumerate}[leftmargin=1.5em]
  \item What did the agent know at each step?
  \item Why did the agent choose the next action?
  \item Which tool calls, observations, and evidence led to the final outcome?
\end{enumerate}

\section{Agent Execution Record}

The AER is the provenance object produced for an agent run. It groups the run
metadata and the ordered sequence of execution steps into a normalized schema.
At the run level, an AER contains identifiers and lifecycle state such as:

\begin{itemize}[leftmargin=1.5em]
  \item \texttt{run\_id}: the unique execution identifier.
  \item \texttt{jira\_ticket\_id}: the originating Jira ticket or task context.
  \item \texttt{agent\_name}: the agent that produced the trace.
  \item \texttt{status}: the execution state, such as running, completed, or failed.
  \item \texttt{started\_at} and \texttt{completed\_at}: timestamps for the run.
  \item \texttt{steps}: the sequence of recorded reasoning and action steps.
\end{itemize}

Each AER step records the local reasoning state. The step schema supports LLM
steps, tool steps, observations, and decisions. A step includes:

\begin{itemize}[leftmargin=1.5em]
  \item \texttt{intent}: what the agent was trying to do.
  \item \texttt{observation}: what the agent observed at that point.
  \item \texttt{inference}: the reasoning or interpretation produced from the observation.
  \item \texttt{context}: the local state visible to the agent.
  \item \texttt{llm}: prompt, system message, model, response, latency, and token count.
  \item \texttt{tool}: tool name, parameters, result, latency, errors, and mock status.
  \item \texttt{evidence}: links to supporting artifacts, verdicts, and confidence.
  \item \texttt{timestamp}: the time at which the step was recorded.
\end{itemize}

\section{Reasoning Provenance Features}

\begin{table}[h]
\centering
\begin{tabularx}{\textwidth}{@{}p{0.28\textwidth}Y@{}}
\toprule
\textbf{Feature} & \textbf{Explainability Value} \\
\midrule
Prompt and response capture &
Shows exactly what the model received and produced at each LLM step. \\
Context snapshots &
Preserves the state available to the agent when it made a decision. \\
Intent, observation, inference fields &
Separates the agent's goal, sensory input, and interpretation so reasoning can
be audited step by step. \\
Tool call provenance &
Records tool parameters, results, latency, errors, and whether a call was
side-effecting or mocked. \\
Evidence chains &
Connects steps to supporting artifacts, verdicts, confidence scores, and
metadata. \\
Timing metadata &
Supports latency analysis across LLM calls, tool calls, and idle time between
steps. \\
Trace hashing &
Provides a stable fingerprint for detecting whether a recorded trajectory has
changed. \\
\bottomrule
\end{tabularx}
\end{table}

\section{Execution DAG}

The visualized DAG presents the AER as an execution graph. Nodes represent
recorded reasoning, decision, LLM, tool, or observation steps. Edges represent
the execution order and dependency between steps. This graph makes the agent's
trajectory readable as a causal path rather than as a flat log.

The DAG view provides:

\begin{itemize}[leftmargin=1.5em]
  \item A run-level view of the full execution trajectory.
  \item Step-level expansion into prompts, responses, tool calls, and evidence.
  \item Visual distinction between LLM reasoning, tool execution, observations,
  errors, and mocked side-effecting calls.
  \item Branch visibility when a replay diverges from the original execution.
  \item A direct path from graph nodes back to the underlying AER fields.
\end{itemize}

\section{Why This Is the Explainability Layer}

gommage provides explainability by sitting between the agent and the systems it
uses. It captures model interaction, tool interaction, context, evidence, and
execution timing without requiring the developer to reconstruct behavior from
unstructured logs. The result is a provenance layer that supports debugging,
audit, replay, and human inspection of agent reasoning.

The AER is the durable representation of that provenance. The DAG is the human
interface for navigating it. Together, they make agent behavior explainable at
the level where failures actually occur: individual reasoning and action steps.

\end{document}
